\documentclass[12pt, a4paper, oneside]{article}

%\renewcommand{\familydefault}{\sfdefault}
%\usepackage{helvet}


% Seitenränder gemäß Richtlinie
\usepackage[top=20mm,bottom=20mm,left=30mm,right=25mm,includehead]{geometry}

\usepackage{csquotes}
\usepackage[round,sort&compress, authoryear]{natbib}
\setcitestyle{authoryear,round,aysep={, },yysep={;},notesep={, }}
\bibliographystyle{plainnat}

%Tweaks für scrbook
\usepackage{scrhack}

\usepackage{tabularx} % in the preamble

%Blindtext
\usepackage{blindtext}

%Erlaubt unter anderem Umbrüche captions
\usepackage{caption}
\captionsetup{font=normalsize}

%Stichwortverzeichnis
\usepackage{imakeidx}

%Kompakte Listen
\usepackage{paralist}

%Zitate besser formatieren und darstellen
\usepackage{epigraph}

%Glossar, Stichwortverzeichnis (Akronyme werden als eigene Liste aufgeführt)
\usepackage[toc, acronym]{glossaries}

%Anpassung von Kopf- und Fußzeile
%beeinflusst die erste Seite des Kapitels
\usepackage[automark,headsepline]{scrlayer-scrpage}
%\input{resources/styles/header_footer}
\setlength{\parindent}{0em}
\setlength{\parskip}{6pt}
\clearpairofpagestyles
\ohead{\normalfont\normalsize\pagemark}
\renewcommand*{\headfont}{\normalfont}
\pagestyle{scrheadings}

%\renewcommand\chapter{\thispagestyle{plain}}

%Auskommentieren für die Verkleinerung des vertikalen Abstandes eines neuen Kapitels
%\renewcommand*{\chapterheadstartvskip}{\vspace*{.25\baselineskip}}
%\renewcommand*{\chapterheadendvskip}{\vspace*{.25\baselineskip}}

%Zeilenabstand 1,5
\usepackage[onehalfspacing]{setspace}
\usepackage{etoolbox}
\AtBeginEnvironment{thebibliography}{\singlespacing}
\AtBeginEnvironment{table}{\singlespacing}
\AtBeginEnvironment{tabular}{\singlespacing}
\AtBeginEnvironment{longtable}{\singlespacing}

%Verbesserte Darstellung der Buchstaben zueinander
\usepackage[stretch=10]{microtype}

%Deutsche Bezeichnungen für angezeigte Namen (z.B. Inhaltsverzeichnis etc.)
\usepackage[ngerman]{babel}

% Hauptschrift gemäß Richtlinie
\usepackage{fontspec}
\setmainfont{Times New Roman}
\setlength{\emergencystretch}{1em}

%Einfachere Zitate
\usepackage{epigraph}

%Unterstützung der H Positionierung (keine automatische Verschiebung eingefügter Elemente)
\usepackage{float}

%Erlaubt Umbrüche innerhalb von Tabellen
\usepackage{tabularx}

%Erlaubt Seitenumbrüche innerhalb von Tabellen
\usepackage{longtable}

%Erlaubt die Darstellung von Sourcecode mit Highlighting
\usepackage{listings}

%Definition eigener Farben bei Nutzung eines selbst vergebene Namens
\usepackage[table,xcdraw]{xcolor}

%Vektorgrafiken
\usepackage{tikz}

%Grafiken (wie jpg, png, etc.)
\usepackage{graphicx}

%Grafiken von Text umlaufen lassen
\usepackage{wrapfig}

%Erlaubt das Linebreaking von URLs an beliebiger alphanummerischen Stelle 
\usepackage{xurl}
%Ermöglicht Verknüpfungen innerhalb des Dokumentes (e.g. for PDF), Links werden durch "hidelink" nicht explizit hervorgehoben
\usepackage[hidelinks,german]{hyperref}
%Aktiviert das Linebreaking von URLs. Ist extrem wichtig bei sher langen Links, dass diese nicht über den Seitenrand laufen
\hypersetup{breaklinks=true}
